
\documentclass[preprint,12pt,authoryear]{elsarticle}

\usepackage{amsmath}
\usepackage{amssymb}
\usepackage{booktabs}
\usepackage{array}
%% needspace stops a caption or a part heading from being printed when too
%% little of the page remains for the table that must follow it. It changes
%% page breaking only, never content.
\usepackage{needspace}
\usepackage{longtable}
\newcolumntype{L}[1]{>{\raggedright\arraybackslash}p{#1}}

\newcommand{\us}{\_\allowbreak}
\newcommand{\tablehead}[2]{%
  \par\addvspace{\bigskipamount}%
  \needspace{12\baselineskip}%
  \refstepcounter{table}\label{#1}%
  {\footnotesize\noindent\textbf{Table \thetable.}~#2\par}%
  \nopagebreak\addvspace{4pt}\nopagebreak}
\newcommand{\tabletail}{\par\addvspace{\bigskipamount}}

\renewcommand{\thetable}{S\arabic{table}}
\renewcommand{\thesection}{S\arabic{section}}

\tolerance=1500
\setlength{\emergencystretch}{3em}

\journal{Smart Agricultural Technology}

\begin{document}

\begin{center}
{\Large\bfseries Supplementary material}\\[4pt]
{\large Offline On-Device Two-Stage Crop Disease Detection for Smartphone
Deployment in Ghana}\\[6pt]
\end{center}

\noindent This document provides supporting evidence for the main manuscript.
The emulator and simulator results in Section~\ref{sup:harness} were collected
only to validate the benchmark instrumentation and application behaviour.
They are not measurements of physical-device performance and are not combined
with the iPhone or Xiaomi results reported in the main manuscript.

\section{Physical-device benchmark evidence}
\label{sup:physical}

The physical-device results are presented in the main manuscript and are not
duplicated here. They were derived from one complete session on an
iPhone~15 Pro Max (20 November 2026, release build) and one complete session
on a Xiaomi Redmi Note~14 Pro+~5G (14 September 2026). Earlier iPhone sessions
(three rejection-path sessions of 14 September 2026 and a debug-build session
of 19 November 2026) are archived but not reported.

To confirm that the phones compute what was evaluated on the desktop, both ran
30 checksum-locked images (24 held-out photographs, four non-crop photographs
and two EXIF-rotated copies) through the production pipeline on 20 November
2026. Both returned the desktop classifier label for every accepted image and
rejected all four non-crop photographs; one image whose confidence lay at the
0.55 floor was routed differently. Against the desktop deployed detector path
(one-to-many head, $640\times640$ letterbox, class-wise suppression at IoU
0.10), Android detections matched exactly (189 of 189 boxes) and the iPhone
matched 182 of 189 (mean score difference 0.017). The full comparison is in
the \texttt{output\_agreement} folder of the evidence package. All benchmark
records are archived in the evidence package.

Each value presented in the physical-device tables was independently
recalculated from the corresponding benchmark record. The verification
included latency summaries, routing outcomes, failure counts, thermal
conditions, battery state and model-version consistency.

\section{Emulator and simulator validation}
\label{sup:harness}

\tablehead{tab:S1}{Android-emulator and iOS-simulator benchmark validation.
These results confirm that the shared Kotlin Multiplatform instrumentation
operates correctly on both platforms and produces complete, internally
consistent measurements. They are not physical-device latency results and are
not combined with the physical-device measurements in the main manuscript.
All latency values are reported in milliseconds.}

\needspace{20\baselineskip}
\noindent\textbf{Part A. Android-emulator quick benchmark}

The quick benchmark was repeated in two sessions on an Android phone emulator.
Each condition included 10 warm-up iterations followed by 50 measured
iterations.

\begin{center}
\scriptsize
\setlength{\tabcolsep}{4pt}
\begin{tabular}{c l r r r r r r}
\toprule
Session & Stage & Mean & Min & Max & p50 & p95 & p99 \\
\midrule
1 & classifier                    & 17.994 & 16.990 & 40.342 & 17.548 & 18.159 & 19.048 \\
1 & detector @$1920\times1080$    & 43.739 & 42.307 & 47.853 & 43.649 & 44.964 & 46.141 \\
1 & detector @$1280\times960$     & 43.637 & 42.078 & 52.034 & 43.229 & 46.715 & 47.004 \\
1 & detector @$640\times480$      & 43.743 & 41.242 & 60.960 & 42.607 & 49.868 & 51.325 \\
2 & classifier                    & 17.568 & 16.812 & 18.796 & 17.496 & 18.183 & 18.329 \\
2 & detector @$1920\times1080$    & 44.189 & 42.932 & 47.879 & 44.100 & 45.396 & 45.954 \\
2 & detector @$1280\times960$     & 44.166 & 42.705 & 59.619 & 43.752 & 45.450 & 46.112 \\
2 & detector @$640\times480$      & 42.432 & 41.217 & 46.410 & 42.320 & 43.344 & 44.775 \\
\bottomrule
\end{tabular}
\end{center}

{\footnotesize
The two sessions are reported separately because their upper-tail measurements
differed. For example, the 95th-percentile detector latency at
$640\times480$ was 49.868~ms in the first session and 43.344~ms in the
second, despite similar median values. The emulator also reported software
rendering during startup because of memory pressure on the host computer.
\par}

\medskip

\needspace{20\baselineskip}
\noindent\textbf{Part B. Extended benchmark comparison}

The extended benchmark was run on an Android emulator and an iOS simulator.
Each platform completed 10 warm-up iterations followed by 100 measured
iterations.

\begin{center}
\footnotesize
\setlength{\tabcolsep}{4pt}
\begin{tabular}{p{168pt} r r}
\toprule
Metric & Android emulator & iOS simulator \\
\midrule
Classifier initialisation, mean ($n=5$)       & 0.322 ms & 0.354 ms \\
Detector initialisation, mean ($n=5$)         & 0.220 ms & 0.311 ms \\
Image decoding, mean                          & 0.211 ms & 0.129 ms \\
Orientation correction, mean                  & 0.055 ms & 0.001 ms \\
Preprocessing, mean                           & 0.449 ms & 7.154 ms \\
Classifier inference, mean                    & 17.040 ms & 10.841 ms \\
Routing decision, mean                        & 0.000 ms & 0.000 ms \\
Detector inference @$1920\times1080$, mean    & 38.199 ms & 19.248 ms \\
Detector inference @$1280\times960$, mean     & 38.274 ms & 19.343 ms \\
Detector inference @$640\times480$, mean      & 39.300 ms & 19.007 ms \\
Output-buffer transfer @$1920\times1080$, mean & 0.000 ms & 1.386 ms \\
End-to-end rejection path, mean               & 18.933 ms & 27.803 ms \\
End-to-end rejection path, p95                & 19.130 ms & 32.213 ms \\
Detector executed / skipped                   & 0 / 100 & 0 / 100 \\
Successful / failed iterations                & 100 / 0 & 100 / 0 \\
\bottomrule
\end{tabular}
\end{center}

{\footnotesize
\textit{Notes.}
(1)~The Android output-buffer transfer is reported as 0.000~ms because
Android already exposes detector output as a directly usable numerical array.
The iOS implementation requires a separate transfer from the native inference
layer into shared application memory. This difference reflects the platform
implementations and is not a missing Android measurement.

(2)~The two columns should not be interpreted as a comparison of Android and
iOS hardware. Both environments ran on the same desktop computer using
different emulation and virtualisation systems, resource allocations and
rendering paths. They are shown together only to demonstrate that the
instrumentation operates correctly on both platforms.

(3)~The values do not represent the performance of the physical Xiaomi or
iPhone devices.

(4)~Both runs exercised only the out-of-distribution rejection path because
the synthetic benchmark image was rejected by the classifier. Accepted-crop
and crop-mismatch measurements are unavailable rather than zero.

(5)~A separate six-case Android-emulator evaluation used real photographs and
the application interface to verify accepted crops, out-of-distribution
rejection, crop mismatch and image-orientation handling. That behavioural
test is distinct from the repeated synthetic-image benchmark summarised here.
\par}

\tabletail

\tablehead{tab:S2}{Simulator-only memory evidence. These measurements are
reported separately so that simulator results are not presented alongside
physical-device memory measurements.}

\needspace{20\baselineskip}
\noindent\textbf{Part A. Memory before and after the measured loop}

\begin{center}
\footnotesize
\setlength{\tabcolsep}{5pt}
\begin{tabular}{p{150pt} r r}
\toprule
Environment & Before measured loop & After measured loop \\
\midrule
iOS simulator & 365{,}024 KB (356.5 MB) & 365{,}376 KB (356.8 MB) \\
\bottomrule
\end{tabular}
\end{center}

\medskip

\needspace{20\baselineskip}
\noindent\textbf{Part B. Extended simulator memory test}

The extended memory test included 20 warm-up detector calls followed by 500
measured calls.

\begin{center}
\footnotesize
\setlength{\tabcolsep}{8pt}
\begin{tabular}{l r}
\toprule
Measurement point & Resident memory \\
\midrule
Start              & 385.7 MB \\
After warm-up      & 441.8 MB \\
Call 50            & 441.4 MB \\
Call 100           & 395.4 MB \\
Call 200           & 395.7 MB \\
Call 300           & 395.7 MB \\
Call 400           & 368.2 MB \\
Call 500           & 367.9 MB \\
Peak               & 443.5 MB \\
Final              & 367.9 MB \\
\bottomrule
\end{tabular}
\end{center}

{\footnotesize
Resident memory did not show continuous growth during the 500 measured calls.
After increasing during warm-up, it stabilised and later declined, ending
73.9~MB below the post-warm-up value. This test used the corrected contiguous
output-buffer transfer. It was designed to examine sustained simulator memory
behaviour and should not be interpreted as physical-device evidence.
\par}

\tabletail

\section{Benchmark conditions}
\label{sup:conditions}

\begin{enumerate}
  \item \textbf{iPhone~15 Pro Max.} One complete session was recorded on
  iOS~27.2.1 while the device was charging at 95\% battery from a release build. iOS reported the
  thermal state as fair. All five routing paths were measured with the same
  checksum-locked photographs as on the Xiaomi device.

  \item \textbf{Xiaomi Redmi Note~14 Pro+~5G.} One complete session was
  recorded while the device was charging at 100\% battery, and Android
  reported its lowest thermal category. Charging is
  reported as part of the measurement condition; no claim is made that it
  improved or reduced performance.

  \item \textbf{Iteration counts.} Each extended benchmark used 10 warm-up
  iterations followed by 100 measured iterations for every available routing
  outcome.

  \item \textbf{Inference backend.} Both applications used ExecuTorch with
  the XNNPACK CPU backend. Inspection of the Android runtime found no active
  path for NNAPI, Vulkan, QNN, Hexagon or the Snapdragon neural processing
  unit. The reported measurements therefore represent CPU inference rather
  than accelerator-assisted execution.
\end{enumerate}

\section{Build and model provenance}
\label{sup:provenance}

Both reported sessions were produced from clean, committed release builds
(iPhone revision \texttt{7859f987c1c0}, Xiaomi \texttt{f5193cb92d24}). The two
revisions differ in benchmark instrumentation, an Xcode build-script check and
the output-agreement mode; the inference, routing and output-transfer code is
identical. The platforms are
nevertheless presented as separate evaluations rather than as a controlled
comparison.

Builds containing uncommitted tracked changes are identified automatically by
a suffix on the recorded revision. Neither reported session carries it.

Both devices used the same four-class EfficientNet-B2 classifier and YOLO26n
detector. The classifier was approximately 30.8~MB and the detector
approximately 9.8~MB. Model checksums were compared across the
physical-device sessions and confirmed that the same model versions were used.

\section{Routing outcomes}
\label{sup:routing}

Five mutually exclusive routing outcomes were evaluated:

\begin{enumerate}
  \item \textbf{Accepted Corn, Pepper or Tomato.} The classifier agrees with
  the crop selected by the user. The detector runs, and its predictions are
  restricted to the disease classes associated with that crop.

  \item \textbf{Out-of-distribution rejection.} The classifier identifies the
  image as Other or does not reach the required confidence threshold. The
  detector is skipped.

  \item \textbf{Crop-mismatch rejection.} The classifier recognises a
  supported crop but disagrees with the user's selection. The detector is
  skipped.
\end{enumerate}

Detector-invocation counters recorded whether detection ran or was skipped for
every measured image. This made the application's fail-closed behaviour
directly verifiable.

\section{Measurement limitations}
\label{sup:limitations}

\begin{enumerate}
  \item \textbf{Timing boundary.} Latency was measured inside the application
  using monotonic system clocks. The measurements include only application
  processing and inference; device-control commands, interface automation and
  communication with the development computer are excluded.

  \item \textbf{Initialisation is not complete cold-start latency.}
  ExecuTorch uses memory mapping and delays part of the model-weight loading
  until the first inference. The recorded initialisation values therefore do
  not represent the complete time required to prepare a model for its first
  prediction and are excluded from the main latency comparisons.

  \item \textbf{Xiaomi memory change is unavailable.} The Xiaomi session's
  before-and-after memory samples were captured consecutively at the end of
  the benchmark rather than on opposite sides of the measured loop. The
  resulting difference was not meaningful and was excluded from all tables
  and claims. The instrumentation was corrected after this session.

  \item \textbf{Platform resource measures differ.} Processor utilisation is
  measured differently by Android and iOS and cannot be compared directly.
  Installed-application size is also represented differently: Android uses
  the application package size, while iOS uses the sum of files in the
  application bundle.

  \item \textbf{Limited device and session coverage.} The physical-device
  evidence consists of one session on one iPhone and one session on one
  Xiaomi device. These results should not be generalised to other units,
  device models or the broader Android and iOS populations.

  \item \textbf{Builds differ by revision.} Both sessions used release builds,
  but from different commits; the output-agreement run on Android used a
  debuggable build of the iPhone's commit.
\end{enumerate}

\section{Evidence traceability}
\label{sup:mapping}

The numerical results were traced to the following retained evidence:

\begin{enumerate}
  \item Detector results were derived from the archived evaluation summaries
  and environment records for YOLO26n, YOLO26s, YOLO26m and RT-DETR-L.

  \item Classifier results were derived from the archived three-class and
  four-class evaluation summaries.

  \item Composed-pipeline results were derived from the archived evaluation
  summaries and routing records for YOLO26n and RT-DETR-L under both
  classifier configurations.

  \item Physical-device latency and routing results were derived from the
  retained iPhone benchmark record of 20 November 2026 and the single
  retained Xiaomi benchmark record.

  \item Android behavioural-acceptance results were derived from the
  six-case application log and detector-invocation records.

  \item The iOS memory trace was derived from the phase-boundary console log
  and memory samples of the iPhone session of 20 November 2026. The
  pre-correction termination point comes from the session record of
  14 September 2026; that process left no export.

  \item Supplementary Table~S1 was derived from two Android-emulator quick
  benchmarks and one extended benchmark from each simulated platform.

  \item Supplementary Table~S2 was derived from the iOS-simulator extended
  benchmark and the 500-call memory test.

  \item Dataset-overlap findings were derived from the retained average-hash
  audits comparing the test images with the training and validation images.
\end{enumerate}

\section{Optional future measurements}
\label{sup:future}

The following measurements may strengthen the evidence but are not required
to interpret the results currently reported:

\begin{enumerate}
  \item Collect unplugged sessions or sessions beginning from a cooler device
  state to examine the sensitivity of latency to charging and accumulated
  heat.

  \item Repeat both device evaluations from the same committed source version
  to support a more controlled cross-platform comparison.

  \item Repeat the complete Xiaomi benchmark to characterise run-to-run and
  thermal variation.

  \item Complete the cross-platform output-agreement evaluation to compare
  desktop, Android and iOS classification, routing and detection outputs,
  including image-orientation handling.
\end{enumerate}


\section{Dataset audit and rebuilt split}
\label{sup:audit}

Two views of the same photograph collection exist in the modelling repository.
The \emph{exported view} is the Roboflow export used to train and evaluate every finalised model: 5{,}239 photographs split into 3{,}199 train, 1{,}024 validation and 1{,}016 test images, carrying 24{,}530, 8{,}247 and 8{,}241 boxes, at a 640-pixel long side.
The \emph{consolidated view} holds 58{,}361 files under UUID names at $640 \times 640$ pixels, with exactly one annotation row per file.
Grouping its files by SHA-256 recovers 7{,}004 unique photographs with 8.3 boxes each on average: a repository refactor had flattened a many-boxes-per-photograph dataset into one file per box.

\paragraph{Internal duplication.}
The consolidated view's declared test split of 11{,}672 rows resolved to 4{,}723 unique photographs, of which 4{,}569 (96.7\%) were byte-identical to a file in its train or validation directories.
We therefore rebuilt a photograph-level split.
Photographs, not rows, were assigned to train, validation and test by rarest-class-first greedy stratified assignment under a fixed seed, targeting 70:15:15, with a floor rule that prevents any class's validation or test allocation from falling to zero or one box.
A naive random or crop-stratified assignment had left the rarest class, Pepper Early Blight, with a single test box.
Cross-split SHA-256 overlap in the rebuilt split was re-verified at zero in all three pairwise directions.

\paragraph{Label consistency.}
Conflicting Pepper and Tomato crop labels were found on 215 of the 7{,}004 photographs (3.1\%).
After manual verification, 212 were relabelled as Tomato and three that could not be assigned confidently were excluded.
The rebuilt split holds 7{,}001 photographs (4{,}902 train, 1{,}052 validation, 1{,}047 test) with 40{,}676 training and 8{,}698 test boxes.
Its training partition has a 77.6-fold box imbalance (6{,}595 Corn Cercospora Leaf Spot to 85 Pepper Early Blight) and its test partition a 43.4-fold imbalance.
No reported model was trained on this split.

\paragraph{Cross-view overlap.}
Because the finalised models were trained on the exported view, we compared the consolidated test partition with the exported train and validation images.
Of the consolidated test photographs, 50.4\% matched exactly and 60.5\% within an average-hash Hamming distance of 5; manual inspection of flagged pairs confirmed re-exported copies of the same photographs.
Results on the consolidated test partition were retained for audit only and are not reported in the main text.
Against the consolidated \emph{training} partition, on which the earlier torchvision detectors were trained, 88.0\% of the held-out test photographs matched exactly and 99.0\% within 5 bits, so those runs are not reported; the Faster R-CNN in the main text was retrained on the exported view with complete annotations.

\paragraph{Taxonomy and identifiers.}
The 23 classes were verified as consistent across the two label maps, the detector dataset configuration and the dataset documentation.
YOLO-format labels use identifiers 0 to 22 with no background class; the region-proposal label map uses 1 to 23 with 0 reserved for background.

\section{Per-class held-out detection results}
\label{sup:perclass}

\tablehead{tab:S3}{Per-class AP@0.50 on the held-out test split (1{,}016 images, 8{,}241 boxes), with recall at confidence 0.50 for YOLO26n and RT-DETR-L. Classes with fewer than 100 boxes are indicative only.}
\begin{center}
\footnotesize
\setlength{\tabcolsep}{4pt}
\begin{tabular}{l r r r r r r r}
\toprule
 & & \multicolumn{4}{c}{AP@0.50} & \multicolumn{2}{c}{Recall @ 0.50} \\
\cmidrule(lr){3-6}\cmidrule(lr){7-8}
Class & Boxes & 26n & 26s & 26m & RT-DETR-L & 26n & RT-DETR-L \\
\midrule
Corn Cercospora Leaf Spot & 1{,}338 & 0.3430 & 0.3759 & 0.3740 & 0.3872 & 0.1151 & 0.4723 \\
Corn Common Rust & 310 & 0.2915 & 0.2492 & 0.3072 & 0.3111 & 0.1194 & 0.3613 \\
Corn Healthy & 290 & 0.2824 & 0.2932 & 0.2750 & 0.3098 & 0.1379 & 0.3931 \\
Corn Northern Leaf Blight & 33 & 0.2013 & 0.2006 & 0.3293 & 0.4402 & 0.0303 & 0.4242 \\
Corn Streak & 721 & 0.3822 & 0.3877 & 0.3753 & 0.4367 & 0.0749 & 0.5562 \\
Pepper Bacterial Spot & 327 & 0.1995 & 0.2140 & 0.2161 & 0.2404 & 0.0245 & 0.2569 \\
Pepper Cercospora & 89 & 0.2528 & 0.2884 & 0.2963 & 0.3146 & 0.0337 & 0.3146 \\
Pepper Early Blight & 8 & 0.0000 & 0.0000 & 0.0054 & 0.0000 & 0.0000 & 0.0000 \\
Pepper Fusarium & 94 & 0.4887 & 0.4877 & 0.4726 & 0.6072 & 0.4468 & 0.6809 \\
Pepper Healthy & 141 & 0.5859 & 0.6370 & 0.5902 & 0.5920 & 0.5248 & 0.7234 \\
Pepper Late Blight & 127 & 0.0617 & 0.0580 & 0.0605 & 0.1017 & 0.0079 & 0.1102 \\
Pepper Leaf Blight & 87 & 0.3341 & 0.4320 & 0.3597 & 0.4582 & 0.1264 & 0.4253 \\
Pepper Leaf Curl & 331 & 0.2426 & 0.2583 & 0.2767 & 0.2597 & 0.1631 & 0.2991 \\
Pepper Leaf Mosaic & 526 & 0.0444 & 0.0562 & 0.0343 & 0.0594 & 0.0019 & 0.1122 \\
Pepper Septoria & 483 & 0.1444 & 0.1622 & 0.1841 & 0.2198 & 0.0269 & 0.1263 \\
Tomato Bacterial Spot & 84 & 0.3081 & 0.3044 & 0.2737 & 0.4319 & 0.1071 & 0.4524 \\
Tomato Early Blight & 465 & 0.3820 & 0.4384 & 0.3990 & 0.4470 & 0.2043 & 0.4710 \\
Tomato Fusarium & 207 & 0.3432 & 0.3323 & 0.3867 & 0.3445 & 0.2899 & 0.4831 \\
Tomato Healthy & 486 & 0.6345 & 0.6622 & 0.6724 & 0.6871 & 0.4630 & 0.7695 \\
Tomato Late Blight & 813 & 0.1586 & 0.1564 & 0.1880 & 0.2195 & 0.0541 & 0.2891 \\
Tomato Leaf Curl & 73 & 0.3093 & 0.3039 & 0.2873 & 0.3578 & 0.1781 & 0.3836 \\
Tomato Mosaic & 25 & 0.3052 & 0.3267 & 0.4124 & 0.4489 & 0.1200 & 0.3600 \\
Tomato Septoria & 1{,}183 & 0.2012 & 0.2101 & 0.2271 & 0.2727 & 0.0482 & 0.2857 \\
\bottomrule
\end{tabular}
\end{center}
\tabletail

\section{Full physical-device records}
\label{sup:devicetables}

The main-text latency table summarises these records. The tables below give
the full per-session detail for each phone, the routing counters, and the iOS
memory trace behind the main-text memory figure. Table and section references
inside them refer to this supplement unless marked as main text.

%% Source: paper_evidence/mobile/iphone_15_pro_max_2026-11-19/extended_benchmark_1795105797518.json
%% Rewritten 2026-11-19 (replaces the three rejection-path sessions of 2026-09-14).
\tablehead{tab:S4}{Physical iPhone~15 Pro Max mobile latency by routing path
during one session (19 November 2026) using a clean, committed debug build.
All timing values are reported in milliseconds.}
\begin{center}
\scriptsize
\setlength{\tabcolsep}{3pt}
\begin{tabular}{l c r r r r r r c c}
\toprule
      & Warm/ &      &        &     & Std- &     &     & Det.\ exec./ & Succ./ \\
Path  & meas. & Mean & Median & p95 & dev  & Min & Max & skipped      & fail.  \\
\midrule
Accepted Corn       & 10/100 & 148.19 & 148.09 & 150.30 & 1.38 & 145.72 & 151.73 & 100 / 0 & 100 / 0 \\
Accepted Pepper     & 10/100 & 157.44 & 157.45 & 167.06 & 7.15 & 147.96 & 184.59 & 100 / 0 & 100 / 0 \\
Accepted Tomato     & 10/100 & 172.16 & 172.00 & 189.18 & 7.03 & 160.42 & 193.28 & 100 / 0 & 100 / 0 \\
OOD rejection       & 10/100 & 51.80  & 50.76  & 59.44  & 3.06 & 49.45  & 65.75  & 0 / 100 & 100 / 0 \\
Crop mismatch       & 10/100 & 50.75  & 50.67  & 52.46  & 1.08 & 49.15  & 56.89  & 0 / 100 & 100 / 0 \\
\bottomrule
\end{tabular}
\end{center}

{\footnotesize
\textit{Device and models.} iPhone~15 Pro Max (\texttt{iPhone16,2}), iOS~27.2.1,
8~GB RAM; charging at 45\% battery; thermal state fair. Build revision
\texttt{1ca6261b8df8} (no uncommitted tracked changes). Classifier
30{,}820{,}096~B (SHA-256 \texttt{39e1b63dac8d\ldots}) and YOLO26n detector
9{,}768{,}796~B (\texttt{baa36404dbb6\ldots}), the same files as on the Xiaomi
device. The end-to-end paths used the same checksum-locked photographs as the
Xiaomi session.

\textit{Isolated-stage figures} (separate measured loops on synthetic patterns,
so they need not sum to the end-to-end means): classifier inference mean
20.77~ms; preprocess 8.64~ms; image decode 0.12~ms. Detector inference mean by
input resolution: $1920\times1080$, 50.50~ms; $1280\times960$, 52.81~ms;
$640\times480$, 54.07~ms; post-processing 0.91--0.92~ms and output transfer
1.69--1.76~ms.

\textit{Notes.} (1)~Process-level CPU utilisation was 1.48 (iOS
\texttt{getrusage}), which is not comparable with the Android per-thread
figure. (2)~Sub-millisecond model-initialisation readings are excluded because
the runtime loads model data lazily. (3)~The headless benchmark entry point is
compiled only into debug builds. (4)~Three earlier sessions of 14 September
2026 (iOS~26.6.2) covered the rejection path only, from a build with
uncommitted changes; they are archived in the evidence package but not
reported.
\par}
\tabletail

\tablehead{tab:S5}{Physical Xiaomi Redmi Note~14 Pro+~5G mobile latency by
routing path during one session using a clean release build. All timing values
are reported in milliseconds.}
\begin{center}
\scriptsize
\setlength{\tabcolsep}{3pt}
\begin{tabular}{l c r r r r r r c c}
\toprule
      & Warm/ &      &        &     & Std- &     &     & Det.\ exec./ & Succ./ \\
Path  & meas. & Mean & Median & p95 & dev  & Min & Max & skipped      & fail.  \\
\midrule
Accepted Corn       & 10/100 & 147.73 & 148.58 & 157.59 & 10.26 & 137.33 & 207.18 & 100 / 0 & 100 / 0 \\
Accepted Pepper     & 10/100 & 151.81 & 151.29 & 160.91 & 8.61  & 139.42 & 192.18 & 100 / 0 & 100 / 0 \\
Accepted Tomato     & 10/100 & 150.25 & 150.70 & 158.67 & 5.80  & 139.05 & 163.42 & 100 / 0 & 100 / 0 \\
OOD rejection       & 10/100 & 53.64  & 53.12  & 57.58  & 2.07  & 52.33  & 63.94  & 0 / 100 & 100 / 0 \\
Crop mismatch       & 10/100 & 48.93  & 48.55  & 56.99  & 4.40  & 44.74  & 66.20  & 0 / 100 & 100 / 0 \\
\bottomrule
\end{tabular}
\end{center}

{\footnotesize
The five routing outcomes are defined in
the main-text protocol section.

\textit{Isolated-stage figures} (from a separate measured loop using the same
stage-timed functions, so they need not sum exactly to the end-to-end means;
this is run-to-run jitter, not an inconsistency): classifier inference mean
43.23~ms; image decode mean 0.69~ms; orientation-correction mean 0.079~ms;
preprocess mean 0.70~ms. Detector inference mean by input resolution:
$1920\times1080$, 93.75~ms; $1280\times960$, 92.79~ms; $640\times480$,
94.13~ms. Output-buffer transfer accounted for 0.54--0.55~ms in each case and was included
within detector inference. The Android processor-utilisation estimate was
0.98 and was derived from time spent by one processor thread; it is therefore
not directly comparable with the process-level estimate reported for iOS in
Table~\ref{tab:S4}.

\textit{Device and build.} Testing used a clean release build on a physical
Xiaomi Redmi Note~14 Pro+~5G running Android~16 with HyperOS. The device had a
QTI SM7635 processor, a 64-bit ARM architecture and approximately 11.09~GiB
of memory. The battery was at 100\% and charging throughout, while the device
reported its lowest thermal state. Charging is treated as part of the test
condition, and its effect is not estimated. The same classifier and detector
versions were used on both mobile platforms. The installed Android package
was approximately 381~MB, although this package-size estimate is not directly
comparable with the iOS bundle estimate. Inference used the CPU-based XNNPACK
backend; no mobile graphics, neural-network or vendor NPU backend was present.
The reported measurements must therefore not be described as NPU-accelerated.

\textit{Notes.} (1)~Memory samples from this run are not used for any claim
(see Table~\ref{tab:S7}) because the before and after readings were taken at
the same point rather than on opposite sides of the measured loop. This
sampling defect has since been corrected. (2)~Sub-millisecond
model-initialisation readings are excluded for the same reason as in
Table~\ref{tab:S4}. (3)~This is a single session
($n = 1$); run-to-run and thermal-drift variability on this device, as
observed for the iPhone in Table~\ref{tab:S4}, has not yet been
characterised.
\par}
\tabletail

\tablehead{tab:S6}{Two-stage classifier-gated routing outcomes and
fail-closed verification, reported per device and per routing path. Paths,
devices and measurement types are never combined into a single headline
figure. Part~A is physical-device counter evidence; Part~B is emulator
behavioural-acceptance evidence and is not physical-device evidence.}

\needspace{10\baselineskip}\noindent\textbf{Part A. Physical-device harness-measured routing outcomes}

\smallskip
\noindent\textit{Xiaomi Redmi Note~14 Pro+~5G (one session while charging;
$n = 100$ per path)}

\begin{center}\footnotesize\setlength{\tabcolsep}{5pt}
\begin{tabular}{p{115pt} r r r r r}
\toprule
             &     & Det.  & Det.  & Offline & Offline \\
Routing path & $n$ & exec. & skip. & success & failure \\
\midrule
Accepted Corn        & 100 & 100 & 0   & 100 & 0 \\
Accepted Pepper      & 100 & 100 & 0   & 100 & 0 \\
Accepted Tomato      & 100 & 100 & 0   & 100 & 0 \\
OOD rejection        & 100 & 0   & 100 & 100 & 0 \\
Crop mismatch        & 100 & 0   & 100 & 100 & 0 \\
\bottomrule
\end{tabular}\end{center}

\smallskip
\noindent\textit{iPhone~15 Pro Max (three consecutive physical sessions;
$n = 100$ per run; only the out-of-distribution path was exercised in the
100-iteration harness loop in the current evidence set)}

\begin{center}\footnotesize\setlength{\tabcolsep}{5pt}
\begin{tabular}{l l r r r r r}
\toprule
    & Thermal &     & Det.  & Det.  & Offline & Offline \\
Run & state   & $n$ & exec. & skip. & success & failure \\
\midrule
Run 1 & nominal & 100 & 0 & 100 & 100 & 0 \\
Run 2 & fair    & 100 & 0 & 100 & 100 & 0 \\
Run 3 & fair    & 100 & 0 & 100 & 100 & 0 \\
\bottomrule
\end{tabular}\end{center}

{\footnotesize
Across the five Xiaomi paths and three iPhone runs, all operations completed
without an offline failure. The detector was never invoked for an
out-of-distribution image or a crop mismatch, but it was invoked for every
accepted image. These are observed counts rather than statistically estimated
failure rates.
\par}

\medskip
\needspace{10\baselineskip}\noindent\textbf{Part B. Emulator behavioural-acceptance evidence
(supplementary only; not physical-device evidence; not used to support any
deployment-performance claim)}

\begin{center}\footnotesize\setlength{\tabcolsep}{5pt}
\begin{tabular}{p{110pt} l c c l}
\toprule
     & Classifier      & Detector & Detector & Routed \\
Case & outcome         & executed & calls    & classes \\
\midrule
Supported Corn   & Corn @ 0.924    & yes & 1 & Corn diseases   \\
Supported Pepper & Pepper @ 0.891  & yes & 1 & Pepper diseases \\
Supported Tomato & Tomato @ 0.949  & yes & 1 & Tomato diseases \\
Unsupported cassava leaf & Other @ 0.800 & no  & 0 & none        \\
Crop mismatch (Corn image, Tomato selected) & Corn @ 0.924 & no & 0 & none \\
Rotated Corn image & Corn @ 0.925 & yes & 1 & Corn diseases   \\
\bottomrule
\end{tabular}\end{center}

{\footnotesize
Part~B was conducted on an Android phone emulator. Six cases were passed
through the application's normal user interface, while an independent counter
recorded whether the detector was invoked. A separate 100-iteration automated
test used a synthetic image that the classifier consistently rejected. The
detector was skipped in all 100 iterations, and every iteration completed
without failure. The manually driven cases and the automated test are
reported separately and are not combined into a single count.

\textit{Notes.} (1)~\emph{Fail-closed} means that the detector is not invoked
when the classifier rejects an unsupported image or when its prediction does
not match the crop selected by the user. In both situations, the application
still returns a valid response rather than failing. (2)~No accepted-path or
crop-mismatch-path 100-iteration harness data exists for the iPhone in the
current evidence set; this is recorded as a limitation, not filled with an
estimate. (3)~The Xiaomi evidence in Part~A is from a single pilot session
while the device was plugged in and charging, a disclosed measurement
condition. This may affect the timing figures in Table~\ref{tab:S5}, but not
the observed routing counts reported here. (4)~Part~B rows are
Android-only and were run under software emulation on a desktop host; they
are not physical-device performance evidence for the iPhone, the Xiaomi or any
other handset.
\par}
\tabletail

%% Source: paper_evidence/mobile/iphone_15_pro_max_release_2026-11-20/
%%   console_phase_trace.txt (Part A) and
%%   extended_benchmark_1795167684106.json, memorySamples (Part B).
%% Rewritten 2026-11-20 (Release build); the earlier pre/post series had no surviving source.
\needspace{22\baselineskip}
\tablehead{tab:S7}{iOS resident memory with the contiguous output buffer on
the physical iPhone~15 Pro Max (one session, 20 November 2026, iOS~27.2.1,
committed release build). Values are mebibytes (MiB) of resident memory.}

\nopagebreak\noindent\textbf{Part A. Phase-boundary trace}

\begin{center}\footnotesize\setlength{\tabcolsep}{5pt}
\begin{tabular}{p{190pt} r}
\toprule
Phase boundary & Resident memory \\
\midrule
Start of benchmark                       & 34.1 \\
After cold-load loop                     & 34.7 \\
After models and benchmark images loaded & 105.0 \\
After classifier stage (110 calls)       & 141.6 \\
After detector stage, $1920\times1080$   & 205.7 \\
After detector stage, $1280\times960$    & 205.5 \\
After detector stage, $640\times480$     & 204.7 \\
After end-to-end loop (all five paths)   & 327.0 \\
\bottomrule
\end{tabular}\end{center}

\medskip
\noindent\textbf{Part B. Benchmark memory samples}

\begin{center}\footnotesize\setlength{\tabcolsep}{5pt}
\begin{tabular}{p{190pt} r}
\toprule
Sample & Resident memory \\
\midrule
Before the measured loops (after cold-load timing) & 34.9 \\
After all measured loops, including end-to-end     & 300.7 \\
\bottomrule
\end{tabular}\end{center}

{\footnotesize
\textit{Notes.} (1)~The application printed Part~A at each phase boundary;
the value is the process resident size divided by $1024^2$. Part~B's samples are written to the benchmark export at
different instants from the console lines, so the two parts need not agree.
(2)~Neither part samples inside a timed call, so a short-lived peak during a
single inference is not captured.
(3)~The earlier build passed each of the 226{,}800 detector output values to
shared code as a separate object. According to the session record of
14 September 2026, resident memory grew with every detector call until iOS
terminated the application at about 1{,}600~MiB during the first detector
stage. That process left no export, so no pre-correction trace is reported.
(4)~The correction replaced the boxed values with a contiguous numerical
buffer. Model weights, thresholds, class indexing, coordinate transformations
and non-maximum suppression were unchanged. An output-equivalence check in
the iOS Simulator reproduced all 226{,}800 values bit for bit.
(5)~No phase-boundary trace exists for Android. Android already used a
contiguous numerical representation. The Xiaomi memory difference is
excluded because both samples were captured at the same point.
\par}
\tabletail


\section{Reviewed studies}
\label{sup:literature}

Table~\ref{tab:S8} compares the twenty-two primary studies discussed in the
main-text Related Work section on the dimensions that determine cross-study
comparability.

\tablehead{tab:S8}{Reviewed crop-disease recognition studies compared
against the dimensions that determine cross-study comparability. Twenty-two
primary studies; the four reviews and surveys in the collection are cited in
the narrative but omitted here because they report no primary dataset,
architecture, split or metric. \textit{Not reported} means the field is
absent from every retrievable part of the record; \textit{none reported}
means the paper describes its system in enough detail to establish that no
such mechanism or measurement is present. The metric column gives the single
highest figure each study reports for its own system, with the metric named.
Figures in that column are not comparable across rows.}

%% longtable steps the table counter internally, so \thetable inside the
%% environment is one higher than the number \tablehead printed above. Freeze
%% the printed number here and use it in both continuation headers.
\edef\litcomptablenum{\thetable}

\needspace{10\baselineskip}\noindent\textbf{Part A. Task, data and model}

{\tiny\setlength{\tabcolsep}{3pt}
\setlength{\LTcapwidth}{\textwidth}
\begin{longtable}{L{66pt} L{50pt} L{86pt} L{48pt} L{96pt}}
\toprule
Study & Task & Dataset and imaging condition & Crops and classes & Architecture \\
\midrule
\endfirsthead
\multicolumn{5}{l}{\textit{Table \litcomptablenum{} (Part A), continued}}\\
\toprule
Study & Task & Dataset and imaging condition & Crops and classes & Architecture \\
\midrule
\endhead
Mohanty et al. (2016) & Classification & PlantVillage; controlled & 14 crops, 38 classes & AlexNet, GoogLeNet \\
Chen et al. (2024) & Detection & Own database combining PlantVillage, Kaggle and field-collected images (19{,}989 annotated) & 9 crops, 34 disease classes & YOLOv5 with SE, CARAFE, EIoU, Ghost convolution \\
Rathod et al. (2025) & Classification with binary input gate & Customised PlantVillage with curated and augmented images; controlled base & Not reported & DenseNet201, plus a leaf/non-leaf gate; TensorFlow Lite \\
Krishna et al. (2025) & Classification & PlantDoc (internet-scraped) plus approx.\ 3{,}000 web-sourced images & 27 classes combined (PlantDoc 27, web 16) & EfficientNet-B0, EfficientNet-B3, ResNet50, DenseNet201 \\
Mamba Kabala et al. (2023) & Classification, federated & PlantVillage subsets partitioned across simulated learners; controlled & 4 crops (grape, apple, corn, tomato) & ResNet50, DenseNet121, VGG16, InceptionV3, MobileNetV2, ViT-B/16, ViT-B/32 \\
Feng et al. (2024) & Classification & PlantVillage; cassava (26{,}337 field images, Uganda); tomato leaves (32{,}531) & 38 / 5 / 11 classes & EfficientNet convolution plus group-wise transformer \\
Subhan et al. (2026) & Classification & `New Plant Diseases' (augmented), approx.\ 87{,}000 images & 38 classes & EfficientNetB0 with transfer learning \\
Zainab and Mahum (2025) & Classification of leaf regions & PlantVillage; controlled & Not reported & Abstract: small Inception over leaf regions; body: VGG19, ResNet152, DenseNet169, MobileNetV3 \\
Qushtom et al. (2025) & Classification & Wheat fields, Kotli, Azad Kashmir (prior collection of 1{,}567); 450 images used & Wheat; 2 classes & Enhanced CNN with LIME and SHAP \\
Deputy et al. (2023) & Classification & PlantVillage in colour, greyscale and leaf-segmented renderings; controlled & 38 classes & AlexNet, GoogLeNet, ResNet50, InceptionV3 \\
Nikam (2025) & Classification & PlantVillage subset (Kaggle), with lighting and background augmentation & Not reported & CNN, ResNet50, VGG16, EfficientNet, ViT \\
Wang et al. (2024) & Classification & PlantVillage as stated (54{,}306, 38 classes); paper's rice description inconsistent with it & Rice & DC-GAN plus MDFC-ResNet, with SVD preprocessing \\
Al-Shahari et al. (2025) & Classification & Fruit and leaf disease image sets, 250 samples each & 5 classes per set & Vector median filter, DenseNet201 features, spotted-hyena optimiser, recurrent spiking network \\
Khandagale (2025) & Classification & Cotton leaf, grape, soybean and corn datasets, singly and combined & 4 crops; 15 classes stated (Table 1 lists more) & Residual blocks, attention and lightweight layers \\
Singh et al. (2020) & Classification & 54{,}305-image, 38-class set (cites Mohanty et al.; PlantVillage inferred) & 14 crops, 38 classes & AlexNet, deep CNN, residual network, recurrent CNN, MobileNet \\
Pulicherla et al. (2024) & Classification & Offline-augmented PlantVillage plus public cotton, rice and groundnut sets & Not reported & ResNet9 \\
Reddy et al. (2025) & Classification & Annotated potato leaf images; source not named & Potato & CNN comparison including EfficientNet \\
Fathima et al. (2023) & Classification & Kaggle rice-leaf images plus online and on-site photographs & Rice & VGG16 with transfer learning \\
Jadhav et al. (2023) & Classification & Banana leaf images (web, Kaggle, own capture) & Banana; 4 classes & Image processing plus CNN \\
Chavan et al. (2025) & Classification & Own field images, five farms, Kolhapur & Sugarcane; 10 disease classes & DenseNet and a plain sequential CNN \\
Mathew et al. (2023) & Classification & Not reported & Not reported & Leaf-image classifier plus a microcontroller detector unit \\
S et al. (2026) & Classification & PlantVillage; controlled & Multiple crops; count not reported & MobileNetV2 with Flask delivery \\
\midrule
This study & Detection with crop-aware out-of-distribution gate & RAIL Crop Disease (Ghana) field photographs; uncontrolled & 3 crops, 23 classes & YOLO26n, YOLO26s, YOLO26m, RT-DETR-L; EfficientNet-B2 gate \\
\bottomrule
\end{longtable}\addtocounter{table}{-1}\par}

\medskip
\needspace{10\baselineskip}\noindent\textbf{Part B. Evaluation, gating and deployment}

{\tiny\setlength{\tabcolsep}{2.5pt}
\begin{longtable}{L{58pt} L{66pt} L{60pt} L{50pt} L{54pt} L{56pt}}
\toprule
Study & Split design & Best metric reported for own system & Input gating & Deployment target & On-device measurement \\
\midrule
\endfirsthead
\multicolumn{6}{l}{\textit{Table \litcomptablenum{} (Part B), continued}}\\
\toprule
Study & Split design & Best metric reported for own system & Input gating & Deployment target & On-device measurement \\
\midrule
\endhead
Mohanty et al. (2016) & Leaf-level grouping where mappings exist (41{,}112 of 54{,}306 images); several ratios & 99.35\% accuracy & None reported & Smartphone, as motivation & None reported \\
Chen et al. (2024) & Random division 8:1:1; no leakage audit described & mAP@0.50 92.1\% & None reported & Real-time use, asserted from parameter and computation counts & None reported \\
Rathod et al. (2025) & Validation accuracy; no held-out protocol stated & 96\% classification; 99\% gate & Binary leaf / non-leaf & Android, TensorFlow Lite, operates offline & Average prediction time under 2~s on unnamed Android phones; no memory, CPU or thermal data \\
Krishna et al. (2025) & Within-, cross- and combined-dataset; test sets differ by experiment & 80.19\% (combined training, web test) & None reported & None reported & None reported \\
Mamba Kabala et al. (2023) & Federated partitions; varying learners, rounds and local iterations & ResNet50 99.59\%; ViT variants costlier & None reported & Distributed training, not on-device inference & None reported \\
Feng et al. (2024) & 85:15 train-test (SMOTE and augmentation applied) & 99.8\% PlantVillage; 86.9\% cassava & None reported & None reported & None reported \\
Subhan et al. (2026) & 80:20 train-validation & 99.51\% (split described as train/validation) & None reported & Field monitoring, proposed & None reported \\
Zainab and Mahum (2025) & 80:20 train-test & 99.75\% (listed as training; 98.52\% testing) & None reported; extends the label space instead & None reported & None reported \\
Qushtom et al. (2025) & 70:15:15, 75:15:10 and 80:10:10 & Up to 100\%; binary task, test sets of approx.\ 45--68 images & None reported & None reported & None reported \\
Deputy et al. (2023) & 60:40 to 90:10, crossed with three renderings & 0.999 accuracy; $F_1$ up to 1.0 & None reported & Smartphone, as motivation & None reported \\
Nikam (2025) & Not reported & ViT 98.7\%; VGG16 95.1\% & None reported & None reported & None reported \\
Wang et al. (2024) & Train/validation/test; ratios not given & 95.99\% validation; 88.66\% test & None reported & Agricultural Internet-of-Things & None reported \\
Al-Shahari et al. (2025) & 80:20 and 70:30 & 98.60\% (training portion); 98.40\% test & None reported & IoT infrastructure stated; image-only experiments & None reported \\
Khandagale (2025) & 80:10:10 per class & 99.7\% grape; 95.3\% combined (corn 98\% in body) & None reported & Real-time use, asserted from model size & None reported \\
Singh et al. (2020) & 80:20 train-validation & 99\% (abstract); 100\% on 570 test images; 97.5\% validation & None reported & Advisory tool for farmers and extension workers & None reported \\
Pulicherla et al. (2024) & 80:20 train-validation plus 33 test images & 99.2\% (abstract only) & None reported & Browser-delivered web interface & None reported \\
Reddy et al. (2025) & Not reported & 97\% accuracy, EfficientNet & None reported & Cloud-server inference; mobile/web dashboard & Not on-device; inference time stated without value \\
Fathima et al. (2023) & 70:30 train-test, with a further validation split & 95\% test (40 train, 20 test photographs) & None reported & None reported & None reported \\
Jadhav et al. (2023) & Not reported & 88\% testing accuracy & None reported & None reported & None reported \\
Chavan et al. (2025) & Not reported & 94\% / 75\% (abstract); 89.98\% / 73.28\% test & None reported & None reported & None reported \\
Mathew et al. (2023) & Not reported & 70\% test accuracy & None reported & Microcontroller unit writing to an online database & None reported \\
S et al. (2026) & Not reported & Approx.\ 94\% accuracy & None reported & Browser-delivered web platform & None reported \\
\midrule
This study & Inherited export split, audited for near-duplicates after training; residual overlap and its effect reported & mAP@0.50 0.3455, RT-DETR-L, 23-class field test split & Crop-aware out-of-distribution gate, with its end-to-end cost reported & Android and iOS phones, fully offline & Per-path latency, memory, processor utilisation, thermal state and charging condition \\
\bottomrule
\end{longtable}\addtocounter{table}{-1}\par}

{\footnotesize
\textit{Notes.} (1)~The metric column is deliberately heterogeneous and is
not a ranking. All twenty-two primary studies report a metric: twenty-one
report classification accuracy, mostly on controlled or partly controlled
imagery, and one reports detection mAP. Several headline figures are
training, validation or abstract-only values, marked as such. (2)~Only
Mohanty et al.\ describe subject-level (leaf) grouping, and seven studies give
no split ratio. (3)~One study places a model on a handset and reports a
coarse prediction-time bound only; none reports on-device memory, processor
utilisation, thermal state or power condition. (4)~One study reports an input
gate; it is binary and its effect on end-to-end performance is not
quantified. (5)~Every cell was read from the full text of the corresponding
study, item by item; page references are held in the authors' audit record.
\par}
\tabletail



\end{document}